\documentclass[10pt,a4paper]{amsart}
\usepackage[margin=19mm]{geometry}
\usepackage{amsmath,amssymb,mathtools}
\usepackage[scr=rsfs,cal=euler]{mathalfa}
\usepackage{xfrac}
\usepackage[unicode,hidelinks,bookmarks=false]{hyperref}
\newcommand{\ie}{i.e.\ }
\newcommand{\cf}{cf.\ }
\newcommand{\resp}{respectively\ }
\newcommand{\Subsection}{\subsection}
\providecommand{\nobreakdash}{\nobreak\hbox{-}}
\setlength{\emergencystretch}{2em}
\multlinegap0pt
\usepackage{amssymb}
\usepackage{enumerate}
\usepackage{mathtools}
\usepackage{booktabs}
\usepackage{float}
\usepackage{caption}
\usepackage{xcolor}
\newtheorem{thm}{Theorem}
\newtheorem{cor}{Corollary}
\newcommand{\D}{\mathbb{D}}
\newcommand{\dd}{\partial}
\newcommand{\s}{\mathbb{S}}
\title{On Contact Round Surgeries on $(\mathbb{S}^3, \xi_{st})$ and their Diagrams}
\author{Prerak Deep, Dheeraj Kulkarni}
\date{Source: arXiv:2504.06074v3 (CC BY 4.0)}
\begin{document}
\maketitle

\noindent\emph{Source and licence.} On Contact Round Surgeries on $(\mathbb{S}^3, \xi_{st})$ and their Diagrams. Version arXiv:2504.06074v3 (CC BY 4.0), distributed by arXiv under the Creative Commons Attribution 4.0 International licence, http://creativecommons.org/licenses/by/4.0/, as recorded in the arXiv licence metadata for this exact version and shown on the abstract page https://arxiv.org/abs/2504.06074v3. Full licence notice: \texttt{licenses/arxiv-2504.06074-CC-BY-4.0.txt}.\par\smallskip

\noindent\textit{Editorial scope.} Counted unit: the article's cor cor (source0.tex lines 967--969). The article's complete original proof of it is delivered with it (source0.tex lines 971--1068, one \texttt{proof} environment of 98 lines). No mathematics is written, repaired or re-derived: the statement and the proof are the article's own lines, in the article's own notation, and no retained line is substituted or re-flowed.  Retained uncounted as the source-local context the proof needs: lines 891--904.  One declared adaptation: exactly 3 ancillary strings inside the retained spans are omitted (image-inclusion commands and, where the source repeats a label, the repeated label definitions) and no other line differs from the source; the figure environments, with their placement, captions and labels, are kept, and the package records the omitted strings in fidelity receipts bound to the affected bodies.  No retained line contains a citation, so the wrapper carries no bibliography and no bibliography entry.  The retained lines contain the article's own diagram code, which the wrapper typesets with the standard tikz packages; no image file is delivered and the package declares no asset dependency.  Editorial formatting only, in addition: an AMS \texttt{amsart} preamble that restates the article's own declarations and macros used by the retained lines, namely the environments \texttt{thm}, \texttt{cor} on separate counters, together with the standard packages those lines need.  The retained environments print in this document as Theorem 1, Corollary 1 in order of appearance, whereas the article prints the same items with the numbering of its own full document.  The complete native source of the article as distributed in the arXiv e-print for this exact version is bundled unchanged as \texttt{sources/176189/source0.tex}.
\begin{thm}[Contact Bridge Theorem]\label{thm: ContBridgeThm}
The following two statements establish a bridge between a set of contact $(\pm1)$-surgery diagrams and contact round surgery diagrams of nice contact joint pairs. 
\begin{enumerate}
    \item [(1)] For a contact $(\pm1)$-surgery diagram on a Legendrian link $L_1 \cup \cdots \cup L_n \subset \s^3$ of a contact 3-manifold $(M, \xi)$, there is a round surgery diagram consisting of contact joint pairs $L= \bigcup_{i=1}^{n'} (L_{i1}\cup L_{i2})$ with round 1-surgery coefficients $k\in \mathbb{Z} $ on both components $L_{1j}$ and round 2-surgery coefficient $m \in \{\pm1\}$ on $L_{i2}$. 
    \item [(2)]Given a contact round surgery diagram of nice contact joint pairs $L= \bigcup_{i=1}^{n} (L_{1i}\cup L_{i2})$ satisfying the following conditions: 
    \begin{enumerate}
        \item The round 1-surgery coefficient is $k \in  \mathbb{Z}$ on both components, 
        \item and round 2-surgery coefficient is $m\in \{\pm1\}$ on $L_{i2}$. 
        %\item after removing thickened torus for round 2-surgery on a 2-manifold obtained by round 1-surgery, we get $\s^3\setminus int(N(L))$. We assume tight contact structure on $\s^3 \setminus int(N(L))$ is $\xi_{st}|_{\s^3 \setminus int(N(L))}$. 
    \end{enumerate}
    The Legendrian link $L= \bigcup_{i=1}^{n} (L_{1i}\cup L_{i2})$ determines a contact $(\pm1)$-surgery diagram such that for each $i,j$; $L_{ij}$ has contact surgery coefficient $m$. 
\end{enumerate}
    
\end{thm}

\begin{cor}
    Suppose $L$ denotes a contact round surgery diagram consisting of nice contact joint pairs $\bigcup_{i=1}^{n}(L_{i1}\cup L_{i2})$ with contact round 1-surgery coefficient $n\in \mathbb{Z}$ on $L_{ij}$ and contact round 2-surgery coefficient $-1$ on $L_{i2}$. Then, $L$ describes a symplectically fillable contact 3-manifold.  
\end{cor}

\begin{proof}
    We apply Theorem \ref{thm: ContBridgeThm} on the given Legendrian link $L$. 
    We get a contact surgery diagram consisting of contact $(-1)$-surgery components. 
    Therefore, we obtain a symplectically fillable 3-manifold after performing contact round surgery on the given diagram. 

    In the following, we discuss the construction of a particular symplectic filling of the resulting 3-manifold explicitly.
 First, we give the topological description of the handle attachment corresponding to the surgery diagram given in the hypothesis.
Let $X^4 := \s^3 \times [0,1]$. Then $\dd X^4= \s^3\times \{0,1\}$. 
We consider a topological joint pair $K_1 \cup K_2$ in $\s^3\times \{1\}$ with round $1$-surgery coefficients $n_{1}$ and $n_{2}$, and a round $2$-surgery on $K_{2}$ with an integer round $2$-surgery coefficient $m$.

%In order to understand the effect of round surgery along a joint pair in terms of handle attachments, we recall the (topological) notion of a joint pair along a pair of knots $K_{1} \cup K_{2}$.


In the $4$-dimensional setup, we attach a round $1$-handle $R^{1}$, which is a diffeomorphic copy of $\s^1\times \D^1\times \D^2$, to $X^4$ along the attaching region $\s^{1} \times \s^{0} \times \D^{2}$, as shown in Figure~\ref{fig: R1,2H}. 
In Figure~\ref{fig: R1,2H}, we use bricks to model round handles of indices $1$ and $2$. 
Then, we attach a round $2$-handle $R^{2}$, which is a diffeomorphic copy of $\s^{1} \times \D^{2} \times \D^{1}$, along its attaching region $\s^{1} \times \partial \D^{2} \times \D^{1}$.

\begin{figure}[ht]
    \centering
    
    \caption{The left and right bricks represent schematic diagrams of a round $1$-handle and a round $2$-handle, respectively. The left and right faces of each brick are identified. The top and bottom faces of each brick indicate the attaching regions of the handles, while the front and back faces represent their belt regions.}
    \label{fig: R1,2H}
\end{figure}

We recall that, in a joint pair, we perform the round $2$-surgery on a thickened neighborhood of a torus parallel to $\partial N(K_{2})$, which is equivalent to attaching $R^{2}$ to the belt region of the round $1$-handle. 
Furthermore, we attach it in a specific way to mimic the effect of the joint pair on the boundary. For this, we note the following.

While performing round $2$-surgery as part of a joint pair, the meridional disk of the gluing solid torus maps to a curve isotopic to the canonical longitude on the boundary $\partial N(K_{i})$ for each $i=1,2$.\\
This implies that we attach $\s^{1} \times \D^{2} \times \D^{1}$ to the belt region $\s^{1} \times \D^{1} \times \s^{1} \subset R^{1}$ by a gluing map 
\[
\phi: \s^{1} \times \partial \D^{2} \times \D^{1} \longrightarrow \s^{1} \times \D^{1} \times \s^{1}
\]
defined by
\[
\phi\left(\{p\} \times \partial \D^{2} \times \{t\}\right) = \s^{1} \times \{t\} \times \{p\},
\]
for $p \in \s^{1}$ and $t \in \D^{1}$. 

For the attachment, we refer to Figure~\ref{fig:round2handle}.

\begin{figure}[ht]
\begin{center}

\caption{Round $2$-handle (shown with a red boundary) attached along the belt region of a round $1$-handle (shown with a blue boundary). 
The top and bottom faces of the round $1$-handle are attached to $X^{4}$.}
\label{fig:round2handle}
\end{center}
\end{figure}

After the attachment of $R^{2}$, we get
\[
\begin{aligned}
X^{\prime} &= (X \cup R^{1}) \cup R^{2}, \quad \text{and} \\
\partial X^{\prime} &= \left[\partial X \setminus \left(\s^{1} \times \D^{2} \times \s^{0}\right)\right] 
\cup \left(\D^{2} \times \s^{1} \times \s^{0}\right).
\end{aligned}
\]

Hence,
\[
\begin{aligned}
X^{\prime} 
&= \left\{ X \bigcup_{\s^{1} \times \s^{0} \times \D^{2}} \left(\s^{1} \times \D^{1} \times \D^{2}\right) \right\} 
\bigcup_{\phi} \left(\s^{1} \times \D^{2} \times \D^{1}\right) \\
&= X \bigcup_{\s^{1} \times \s^{0} \times \D^{2}} 
\left\{ \left(\s^{1} \times \D^{1} \times \D^{2}\right) \bigcup_{\phi} \left(\s^{1} \times \D^{2} \times \D^{1}\right) \right\} \\
&= X \bigcup_{\s^{1} \times \s^{0} \times \D^{2}} 
\left\{ \left(\s^{1} \times \D^{1} \times \D^{2}\right) \bigcup_{\phi} \left(\D^{2} \times \D^{1} \times \s^{1}\right) \right\}.
\end{aligned}
\]

By the gluing map $\phi$ defined above, we obtain
\[
\left\{ \s^{1} \times \D^{1} \times \D^{2} \right\} 
\bigcup_{\phi} 
\left\{ \D^{2} \times \D^{1} \times \s^{1} \right\} 
= \D^{2} \times \D^{2} \times \s^{0}.
\]
Thus,
\[
X^{\prime} = X \bigcup_{\s^{1} \times \s^{0} \times \D^{2}} 
\left( \D^{2} \times \D^{2} \times \s^{0} \right),
\]
that is, it is equivalent to attaching two $2$-handles to $X^{4}$, as shown in Figure~\ref{fig: RealizeAsTwo2-handles}.

\begin{figure}[ht]
    \centering
    
    \caption{Each of the top and bottom faces of the diagram can be viewed as a $2$-handle attached along the blue-colored subface. 
The green-colored subfaces represent portions of the boundary of the resulting $4$-manifold.}
    \label{fig: RealizeAsTwo2-handles}
\end{figure}

Now consider $X^{4}$ with a symplectic form given by symplectization of the standard contact 3-sphere. 
By Part 2 of the Contact Bridge Theorem (Theorem~\ref{thm: ContBridgeThm}), the contact joint pair with contact round $1$-surgery coefficient $n$ on the Legedrian 
realization of $K_1$ and $K_2$, and contact round $2$-surgery coefficient $(-1)$ on the Legendrian realization of $K_2$ corresponds to contact $(-1)$-surgeries along each component $K_1$ and $K_2$. This is equivalent to performing Weinstein (symplectic) 2-handle attachment to the symplectic manifold $X^{4}$ to obtain $X'$.
Consequently, the symplectic structure extends naturally to the resulting $4$-manifold $X'$. Hence, the boundary manifold $\dd X'$ together with the induced contact structure is symplectically fillable.
\end{proof}

\end{document}
